\section{Prescribed types and random matchings}\label{sec:types}

We construct the two immersed graphs so that their outgoing-label sets
have odd cross-intersections, the condition needed for cancellation.
The same vertex types also provide a squared-word decay estimate.
Conditioning the matchings on large girth will preserve enough control
to prove a diameter bound and, in \Cref{sec:bounded-pattern}, exclude
the path configurations that could defeat the geometric safeguards.

\subsection{The finite geometry}

Fix
\[
q=128,\qquad v=q^2+q+1=16513,\qquad p=\frac{q+1}{v}.
\]
The projective plane over \(\mathbb F_q\) has \(v\) points and \(v\) lines. Each line has \(q+1\) points, each point lies on \(q+1\) lines, two distinct points lie on one common line, and two distinct lines meet at one point. These facts follow by counting one- and two-dimensional subspaces of \(\mathbb F_q^3\).

Let \(T\) consist of its \(v\) points, called \emph{ordinary letters}, and three further \emph{extra letters}. Pair the extras with three distinct ordinary letters, and pair the remaining ordinary letters arbitrarily in twos. This defines a fixed-point-free involution on \(T\), since \(v-3\) is even. In particular, \(|T|=16516\), and the corresponding rose has \(8258\) edges.

Take disjoint vertex sets \(A,B\), each of size \(n\). For each vertex \(x\), prescribe a set \(S_x\subset T\) of outgoing labels as follows:
\begin{enumerate}
\item On either side the ordinary part of \(S_x\) is a line, with every line represented equally often.
\item Within each line class on \(A\), the extra part is each of the three two-element subsets with proportion \(p/2\) apiece, and is otherwise empty.
\item Within each line class on \(B\), the extra part is all three extras with proportion \(p\), and is otherwise empty.
\end{enumerate}
These are exact empirical proportions. We take \(n\to\infty\) through multiples of \(2v^2\), so all counts are integral. For example, at \(n=2v^2\), each line class on \(A\) has \(129\) vertices of each double-extra type and \(32639\) of the empty type; on \(B\) it has \(258\) of the triple-extra type and \(32768\) of the empty type.

Every letter has frequency \(p\) on each side. Two ordinary line parts meet in either \(1\) or \(q+1=129\) points. An extra part on \(A\) meets an extra part on \(B\) in either \(0\) or \(2\) letters. Thus
\begin{equation}\label{eq:odd-intersections}
 |S_x\cap S_y|\ \text{is odd for every }x\in A,\ y\in B.
\end{equation}
Thus the parity condition is built into the types, independently of
how their edges will be matched. The random choice can now be used
entirely to control the geometry without risking this condition.

Choose two vertices \(x_A\in A\) and \(x_B\in B\) in advance. For \(Y=A,B\), write
\[
Y_t=\{x\in Y:t\in S_x\},\qquad |Y_t|=np.
\]
For each inverse pair, choose a bijection \(Y_t\to Y_{\bar t}\), using its inverse for the opposite direction. A matched pair gives an edge labelled \(t\) in one direction and \(\bar t\) in the other. Do this independently for all inverse pairs and both sides before any conditioning. The resulting graphs are denoted
\[
\Gamma_A,\quad\Gamma_B,\qquad
\Gamma=\Gamma_A\sqcup\Gamma_B.
\]
They immerse into \(F\), since every outgoing letter occurs once at each vertex that carries it. Their degrees lie between \(q+1\) and
\(d_*=q+4=132\).
A prescription for one matching edge can be described using either orientation. Whenever prescriptions are counted, an edge and its reverse describe a single prescription.

\subsection{Turn weights}

For a non-cancelling turn \(t,u\), let \(w(t,u)\) be the frequency of \(u\) among the vertices of \(B\) that contain \(\bar t\). Put \(w(t,\bar t)=0\). The corresponding conditional frequency on \(A\) is at most \(w(t,u)\). Indeed, for distinct \(\bar t,u\), the weights on \(B\) are
\[
\begin{array}{c|ccc}
(\bar t,u)&\text{ordinary, ordinary}&
\text{ordinary, extra}&\text{extra, extra}\\ \hline
w(t,u)&1/(q+1)&p&1.
\end{array}
\]
The mixed entry is symmetric; the extra-extra frequency on \(A\) is \(1/2\). In particular, every positive weight lies in \([1/129,1]\).

For a word \(W=t_1\cdots t_h\), define
\begin{equation}\label{eq:word-weight}
P(W)=\prod_{i=1}^{h-1}w(t_i,t_{i+1}).
\end{equation}
Write \(W^{-1}=\bar t_h\cdots\bar t_1\). Equal letter marginals and symmetry of joint incidence give
\begin{equation}\label{eq:reversal-weight}
w(t,u)=w(\bar u,\bar t),\qquad P(W)=P(W^{-1}).
\end{equation}
For non-cancelling turns this follows by dividing the joint frequency of
\(\bar t,u\) by \(p\); for cancelling turns both weights in the first equality are zero.

\begin{lemma}\label{lem:squared-words}
There is a constant \(\delta>0\) such that, for every sufficiently large \(h\),
\begin{equation}\label{eq:squared-words}
 \sum_{W\in T^h}P(W)^2\le e^{-2\delta h}.
\end{equation}
\end{lemma}
\begin{proof}
Let \(M_{t,u}=w(t,u)^2\), and set \(f(t)=5\) if \(\bar t\) is extra and \(f(t)=1\) otherwise. There are precisely three states of weight five.

If \(\bar t\) is ordinary, the ordinary successors different from \(\bar t\) contribute \((v-1)/(q+1)^2=q/(q+1)\) before the weights \(f\) are inserted. At most three such successors have their weight increased by four. The extra successors all have weight one. Hence
\[
(Mf)(t)\le\frac{q}{q+1}+3p^2+\frac{12}{(q+1)^2}
 =\frac{500731261911}{504183783481}<\frac{149}{150}.
\]
If \(\bar t\) is extra, the other two extras contribute two; the ordinary successors contribute at most \(vp^2+12p^2\). Thus
\[
\frac{(Mf)(t)}{f(t)}
 \le\frac{2+vp^2+12p^2}{5}
 =\frac{820350863}{1363395845}<\frac{149}{150}.
\]
Consequently \(Mf\le(149/150)f\). Since \(\mathbf1\le f\) and
\(\sum_t f(t)=16528<e^{10}\),
\[
\sum_{W\in T^h}P(W)^2
=\mathbf1^{\mathsf T}M^{h-1}\mathbf1
\le16528(149/150)^{h-1}
\le e^{-h/300}\qquad(h\ge3002).
\]
For example, \(\delta=1/600\) is admissible.
\end{proof}

\subsection{Girth conditioning and switching}

Fix a sufficiently small positive constant \(c\), and put
\[
L=\lfloor c\log n\rfloor.
\]
It is enough that
\begin{equation}\label{eq:c-choice}
 c\log(2|T|/p)<1,\qquad 2c\log d_*<1;
\end{equation}
one may take \(c=1/100\). We work with the uniform tuple of matchings conditioned on \(\girth(\Gamma)\ge L\). The next lemma both justifies this conditioning and gives a bound within the conditioned space.

\begin{lemma}\label{lem:prescriptions}
For all sufficiently large admissible \(n\), there is a tuple of matchings with \(\girth(\Gamma)\ge L\).
In the uniform distribution on such tuples, condition further on any feasible set of \(s\) prescribed matching pairs. The probability of an additional specified pair, not already prescribed in either direction, is at most
\begin{equation}\label{eq:prescription-bound}
 \frac{1}{np-s-r_n},
 \qquad r_n=O(d_*^{2L+2})=o(n),
\end{equation}
whenever the denominator is positive.
The error \(r_n\) can be chosen uniformly over all prescriptions.
\end{lemma}
\begin{proof}
Before girth conditioning, a feasible simple cycle of length \(k<L\) has at most \(2n^k|T|^k\) vertex-and-letter descriptions. Its distinct edges require \(k\) distinct matching prescriptions, with probability at most \((np-k)^{-k}\). Thus the expected number of edges lying on cycles shorter than \(L\) is at most
\[
2\sum_{k<L}k\,n^k|T|^k(np-k)^{-k}
 \le2\sum_{k<L}k(2|T|/p)^k=o(n)
\]
for sufficiently large \(n\), by \eqref{eq:c-choice}. Choose a tuple with \(o(n)\) such \emph{bad edges}.

Let \(xy\) be a bad edge, oriented with label \(t\). Among the \(np\) edges for this inverse pair on the same side, choose an edge \(uz\), also oriented with label \(t\), that is not bad and whose endpoint set is at distance greater than \(2L\) from \(\{x,y\}\). This is possible: bad edges exclude \(o(n)\) choices, and the maximum-degree ball bound excludes another \(O(d_*^{2L+2})=o(n)\).

Transpose the two matches, replacing \(xy,uz\) by \(xz,uy\). A short cycle using only one new edge would give an old path between the distant endpoint sets. If a short cycle uses both, deleting them leaves old paths pairing the endpoints as either \((x,u),(y,z)\) or \((x,y),(u,z)\). The first case again joins the distant sets. In the second, the \(u\)-to-\(z\) path together with \(uz\) is an old short cycle, contradicting the choice of \(uz\). This includes \(x=y\), when the \(x\)-to-\(y\) path may have length zero. No short cycle is created. The chosen bad edge disappears, and the other removed edge was not bad, so the number of bad edges strictly decreases. Repetition produces a tuple of girth at least \(L\).

Now restrict to any feasible girth-conditioned set with \(s\) previous prescriptions, and consider an additional compatible pair \(xy\). If it is incompatible, its probability is zero. From every tuple containing \(xy\), perform the same transposition with any sufficiently distant edge \(uz\) of the same inverse pair and side that is not previously prescribed. All edges are now nonbad. At least \(np-s-r_n\) choices are available, with one uniform ball-size error \(r_n\). Each output preserves girth and all prior prescriptions.

Given an output and the pair \(xy\) to restore, the inverse operation is unique: the current match of \(x\) is \(z\), and the current preimage of \(y\) is \(u\). Thus the map from a tuple containing \(xy\), together with an admissible transposition, to its output tuple is injective. If \(\Omega_s\) is the conditioned set and \(\Omega_{s,xy}\) is its subset containing \(xy\), then
\[
|\Omega_{s,xy}|(np-s-r_n)\le|\Omega_s|.
\]
This proves \eqref{eq:prescription-bound}.
\end{proof}

In particular, any feasible collection of \(E=O(L)\) distinct edge prescriptions has conditional probability at most
\begin{equation}\label{eq:few-edge-bound}
 (np)^{-E}\exp(o(L)).
\end{equation}
This follows by applying \Cref{lem:prescriptions} successively. Uniformity of its error is important; no independence after conditioning is asserted.

\begin{lemma}\label{lem:diameter}
With conditional probability tending to one, every connected component of \(\Gamma\) has diameter at most \(D_0L\), for a fixed constant \(D_0\).
\end{lemma}
\begin{proof}
Fix a small positive \(\rho\). On either side, require every vertex set of size \(1\le k\le\rho n\) to grow to at least \(2k\) vertices on adjoining all its neighbors. A failure supplies a set of \(2k\) vertices containing at least
\[
r=\lceil(q+1)k/2\rceil
\]
distinct edges: all edges incident to the failing \(k\)-set lie in the larger set and are counted at most twice by its degree sum.

There are at most \((C_1n/k)^{2k}\) choices for the larger set. Its possible typed edge prescriptions number \(O(k^2)\); choosing an unordered \(r\)-subset costs at most \((C_1k)^r\). Decrease \(\rho\) so that \(r\le np/3\) throughout the range. For sufficiently large \(n\), \Cref{lem:prescriptions} then bounds the probability of all \(r\) prescriptions by \((2/(np))^r\). Hence the probability of a failure of size \(k\), on either side, is at most a fixed multiple of
\[
(C_1n/k)^{2k}(C_1k)^r(2/(np))^r
 \le \bigl(C_2(k/n)^{62.5}\bigr)^k,
\]
because \(64.5\le r/k\le65\). Choose \(\rho\) still smaller so that \(C_2\rho^{62.5}<1/2\). Summing over \(k\le\sqrt n\) gives \(o(1)\); the sum over larger \(k\le\rho n\) is at most \(n2^{-\sqrt n}\).

On the resulting expansion event, every ball reaches more than \(\rho n\) vertices in \(R=O(\log n)\) steps. Along a geodesic, centers separated by \(2R+1\) have disjoint radius-\(R\) balls. At most \(1/\rho\) such balls fit in one side. This bounds each component diameter by \(O(\log n)=O(L)\), with a fixed constant.
\end{proof}
